\documentclass[francais]{rtxarticle}

\usepackage{listings}
\usepackage{tikz}
\usetikzlibrary{shapes,fit}

\lstset{%
basicstyle=\ttfamily\footnotesize\color{LightGray},
backgroundcolor=\color{black}
}

\tikzset{warrow/.style={->, >=stealth, color=black, ultra thick}}

\tikzset{graybox/.style={draw,rectangle,rounded corners=3pt,very thick,densely dashed,color=gray!75,text=black}}
\tikzset{redbox/.style={draw,rectangle,rounded corners=5pt,ultra thick,color=rathaxesred,text=white}}
\tikzset{redcontainer/.style={draw,rectangle,rounded corners=5pt,ultra thick,color=rathaxesred,text=white,minimum height=2.5cm,minimum width=6cm}}

\title{Premiers Pas}
\author{Louis Opter}

\rtxdoctype{Documentation Utilisateur}
\rtxdocstatus{Version Finale}
\rtxdocversion{0.3}

\rtxdochistory
{
0.1 & 09/04/2011 & Louis Opter & Traduction de la version Anglaise \\
\hline
0.2 & 19/07/2011 & Louis Opter & Ajout d'une vue d'ensemble de la génération d'un pilote \\
\hline
0.3 & 12/01/2012 & Louis Opter & Mise à jour par rapport à la version Anglaise \\
}

\begin{document}

\maketitle

\begin{abstract}
\rtx\ est un Langage de programmation Spécifique à un Domaine (DSL) qui permet
de décrire des pilotes de périphériques. \rtx\ compile vers des modules noyaux
écrits en C pour n'importe quel système d'exploitation défini.

Ce document explique comment installer \rtx\ sous Windows et les systèmes
d'exploitations de type Unix et comment générer votre premier pilote de
périphérique.

Vous trouverez aussi dans ce document des liens vers la documentation complète
et comment compiler la dernière la version de \rtx.
\end{abstract}

\tableofcontents

\pagebreak

\section{Installation}

Un installateur existe pour Windows, sur les autres systèmes d'exploitation
seulement une installation depuis les sources est supportée pour le moment.
Bien sûr, si vous souhaitez développer sur \rtx\ lui même vous pouvez choisir
d'installer \rtx\ depuis les sources (c.à.d sans l'installateur).

\subsection{Pré-requis}

\rtx\ n'a pas besoin d'être installé sur le même système d'exploitation que
celui que vous visez pour la génération de pilotes. En d'autres mots, vous
pouvez installer \rtx\ sur Windows et générer les sources d'un pilote pour Linux
(alors que vous êtes sous Windows).

Nous avons rendu très facile la compilation de fichier sources écrit en \rtx\
depuis des «~Makefile~» générés par CMake. Vous aurez donc besoin d'installer
\href{http://www.cmake.org/}{CMake} (>= 2.6) en même temps que \rtx. CMake
fonctionne sur tous les systèmes d'exploitation supportés par \rtx\ et possède
un \href{http://www.cmake.org/cmake/resources/software.html}{installateur pour
Windows}\footnote{Quand vous installez CMake sur Windows, faites attention à
bien sélectionner l'option « Ajoutez CMake dans le \texttt{PATH} pour tous les
utilisateurs ».}.

Enfin, si vous souhaitez compiler les sources générées par \rtx\ vous aurez
aussi besoin des sources de Linux sous Linux, du Kit de Développement de Pilote
de Périphérique pour Microsoft Windows sous Windows, etc.

\subsection{Avec l'installateur Windows}

L'installateur fournis pour Windows inclut tout ce dont vous avez besoin pour
écrire des pilotes en utilisant \rtx\ : le compilateur et sa documentation.

Téléchargez l'\href{http://rathaxes.googlecode.com/files/rathaxes-latest.exe}{installeur},
exécuter le puis suivez les instructions sur l'écran.

\subsection{Depuis les sources sur Unix}

Si vous n'utilisez pas Windows vous êtes obligés d'installer \rtx\ depuis les
«~sources~».

Téléchargez la dernière \href{http://rathaxes.googlecode.com/files/rathaxes-src-latest.tar.gz}
{archive des sources}. Une fois l'archive extraite, déplacez vous avec «~cd~» à
l'intérieur. Vous pouvez installer \rtx\ avec les lignes de commandes suivantes~:

\begin{lstlisting}
$ mkdir build
$ cd build
$ cmake -DCMAKE_INSTALL_PREFIX=/usr/local/ -DCMAKE_BUILD_TYPE=RELEASE ..
$ sudo make install
\end{lstlisting}

\emph{Vous aurez besoin d'être root pour exécuter «~make install~», cet exemple
utilise «~sudo~» mais adaptez le à votre installation locale}.

\section{Générez votre premier pilote}

Le processus pour générer un pilote est le même pour chaque système
d'exploitation mais le nom de quelques programmes et certain chemins de fichiers
changent; sélectionnez les instructions qui correspondent a votre système.

\subsection{Depuis Windows}

Open the ``\rtx\ Shell''. Then, ``cd'' somewhere you have read-write rights and
generate the \rtx\ sources for our ``LKM'' example:

Ouvrez le « \rtx\ Shell », ensuite utilisez « cd » pour vous déplacer quelque
part où vous avez les droits en lecture et en écriture puis générez les sources
de notre exemple « LKM » :

\begin{lstlisting}
$ mkdir build
$ cd build
$ cmake "-GNMake Makefiles" %RATHAXES_PATH%/share/doc/samples/LKM/
$ nmake lkm_src
\end{lstlisting}

\subsection{Depuis n'importe quel Unix}

Ouvrez votre émulateur de terminal favoris et générez les sources de notre
exemple de pilote « LKM » :

\begin{lstlisting}
$ mkdir build
$ cd build
$ cmake /usr/share/doc/rathaxes/samples/LKM/
$ make lkm_src
\end{lstlisting}

Si vous êtes sous Linux et avez un arbre des sources valide, vous pouvez aussi
compilez le module noyau que vous venez de générer:

\begin{lstlisting}
$ make lkm
\end{lstlisting}

\section{Plongée dans \rtx}

Avant d'aborder en détail le compilateur et le langage \rtx\ voici une vue
d'ensemble du processus de génération d'un pilote~:

\begin{figure}[ht!]
\begin{center}
\begin{tikzpicture}
\node (RTX) at (-5,0) {\includegraphics[height=1.5cm]{icons/rtx_no_stroke}};
\draw (RTX.south) node[below]{\large{\sffamily{Description du pilote}}};

\node (C) at (5,0) {\includegraphics[height=1.5cm]{icons/c_no_stroke}};
\draw (C.south) node [below]{\large{\sffamily{Module noyau}}};

\node[graybox] (AST1) at (-2,0){\texttt{\large{AST}}};
\draw[warrow] (RTX)--(AST1);

\node[graybox] (AST2) at (2,0) {\texttt{\large{AST}}};
\draw[warrow] (AST2)--(C);

\node[redcontainer] (COMPILER) at (0,0) {};
\draw (COMPILER.south) node[below]{\large{\sffamily{Compilateur}}};

\node[graybox] (RTXLINK) at (COMPILER.center) {\texttt{\large{RTX\_Link}}};
\draw[warrow] (AST1)--(RTXLINK);
\draw[warrow] (RTXLINK)--(AST2);

\node (RTI) at (-1.5,2.7) {\includegraphics[height=1.3cm]{icons/rti_no_stroke}};
\draw (RTI.north west) node[above right]{\large{\sffamily{Interfaces}}};
\draw[warrow] (RTI.south)--(RTXLINK);

\node (BLT) at (0.5,2.7) {\includegraphics[height=1.3cm]{icons/blt_no_stroke}};
\draw (BLT.north west) node[above right]{\large{\sffamily{Patrons de code}}};
\draw[warrow] (BLT.south)--(RTXLINK);
\end{tikzpicture}
\caption{Vue d'ensemble de la génération d'un pilote}
\end{center}
\end{figure}

Vous pouvez constater que le compilateur utilise trois types de fichiers. Ces
fichiers ont tous des extensions différente mais sont en réalité écrit dans le
même langage : \rtx. (Comme en C ou, les fichiers \texttt{.c} et \texttt{.h}
utilisent des extensions différentes mais sont écrits dans le même
langage). Ces trois « types » de fichiers ont des rôles différents :
\begin{description}
\item [{\texttt .rtx}] : décrit votre pilote, ainsi que ces propres algorithmes
qui sont indépendant à chaque plateforme ;
\item [{\texttt .rti}] : contiennent des \emph{interfaces} qui décrivent ce qui
doit être implémenté dans les fichiers \texttt{.rtx} et ce qui doit être fournis
par les \emph{patrons} de code enregistrés dans les fichiers \texttt{.blt} ;
\item [{\texttt .blt}] : contiennent des \emph{unités} de code qui peuvent être
spécifiques à un système d'exploitation; ces unités sont sélectionnées grâce aux
informations fournies dans le fichier \texttt{.rtx} et l'introspection du
système.
\end{description}

Un projet \rtx\ est compilé à partir d'un fichier \texttt{.rtx} et d'un ou
plusieurs fichiers \texttt{.rti} et \texttt{.blt}. La documentation complète de
la partie du langage utilisée dans les fichiers \texttt{.rtx} est disponible
ici : \href{http://rathaxes.googlecode.com/files/dsl_frontend_fr.pdf}{dsl\_frontend\_fr.pdf}.
Un autre document couvre les fichiers \texttt{.rti} et \texttt{.blt} :
\href{http://rathaxes.googlecode.com/files/dsl_backend_fr.pdf}{dsl\_backend\_fr.pdf}.

\section{Installer la version de développement}

Vous pouvez installer la dernière version de \rtx\ si vous avez besoin d'avoir
les derniers bogues et fonctionnalités ou si vous souhaitez contribuer au
projet. Ceci implique de cloner la dernière version du projet en utilisant
Mercurial et de le compiler manuellement.

\subsection{Pré-requis}

Pour cloner et compiler le projet vous avez besoin d'installer ou d'avoir les
logiciels suivants :
\begin{itemize}
\item Mercurial >= 1.5 (vous pouvez vérifier votre version avec la commande
      «~hg {-}{-}version~» et utiliser
      \href{http://tortoisehg.bitbucket.org/download/index.html}{TortoiseHg}
      sur Windows) ;
\item Subversion (vous devez installer
      \href{http://www.sliksvn.com/en/download}{Slik SVN} qui contient les
      programmes en lignes de commande).
\end{itemize}

Vous aurez aussi besoin des pré-requis habituel pour \rtx\ (CMake et une chaîne
de compilation). Si vous souhaitez générer la documentation vous aurez besoin
d'une distribution \LaTeX\ (\href{http://www.tug.org/texlive/}{\TeX Live} >= 2009
ou \href{http://www.miktex.org/}{Mik\TeX} sur Windows) et de
\href{http://pages.cs.wisc.edu/~ghost/}{Ghostscript}.

\subsection{Cloner les sources}

Ouvrez un shell (depuis le menu de Visual Studio sous Windows) et cloner le
projet en utilisant :

\begin{lstlisting}
$ hg clone https://rathaxes.googlecode.com/hg/ rathaxes
$ cd rathaxes
\end{lstlisting}

Gardez le shell ouvert, la section suivante explique comment compiler \rtx\ sur
Windows ou Unix.

\subsection{Compiler \rtx}

Si vous utilisez un système d'exploitation de type Unix utilisez les commandes
suivantes :

\begin{lstlisting}
$ mkdir build
$ cd build
$ cmake ..
$ make
\end{lstlisting}

Si vous êtes sous Windows utilisez :

\begin{lstlisting}
$ mkdir build
$ cd build
$ cmake -G "NMake Makefiles" ..
$ nmake
\end{lstlisting}

\rtxmaketitleblock

\end{document}
